\chapter{实验设计与结果占位}

\section{实验目标}

本文实验的目标不是单纯压测系统上限，而是验证应用层故障注入机制能否支持可靠性分析。一次完整实验需要回答四个问题：

\begin{enumerate}[label=(\arabic*)]
  \item 故障是否按照声明在目标位置实际命中。
  \item 系统在故障条件下的端到端成功率、延迟和错误率如何变化。
  \item Agent 或工具是否触发现有 fallback，业务质量是否发生退化。
  \item 实验是否可以通过相同配置、相同负载和相同记录模板复查或复跑。
\end{enumerate}

因此，实验数据应至少包含三层：Locust 的负载与 HTTP 指标、Tempo 的 Trace 与故障命中证据、Prometheus/Grafana 的系统资源和聚合指标。对于 AI 原生链路，还需要额外记录业务质量指标，例如生成的天数、活动数量、Markdown 长度和持久化成功率。

\section{链路 A 四象限实验矩阵}

链路 A 是当前最适合进行第一阶段实验的对象。本文采用四象限设计，把负载压力和故障条件作为两个维度，如表 \ref{tab:quadrants} 所示。

\begin{table}[htbp]
  \centering
  \caption{链路 A 四象限实验矩阵}
  \label{tab:quadrants}
  \begin{tabularx}{\textwidth}{p{2.8cm}p{2.2cm}p{2.2cm}X}
    \toprule
    实验名 & 压力 & 故障 & 主要用途 \\
    \midrule
    baseline-low & 低压 & 无 & 建立健康稳态，排除系统自身异常 \\
    fault-low & 低压 & 有 & 观察故障本身造成的退化 \\
    baseline-high & 高压 & 无 & 观察正常高并发下容量和瓶颈 \\
    fault-high & 高压 & 有 & 观察高压与故障叠加是否导致级联失败 \\
    \bottomrule
  \end{tabularx}
\end{table}

低压和高压的具体用户数、spawn rate 与持续时间应根据本机和 LLM 网关能力校准。实验记录中必须写明实际参数，不能只写“低压”或“高压”。如果本地机器无法达到文档建议的高压范围，应降低负载并说明原因，保证四象限结构不变。

\section{首批故障场景}

第一阶段建议选择链路 A 中最能体现系统退化差异的代表故障，如表 \ref{tab:first-faults} 所示。

\begin{table}[htbp]
  \centering
  \caption{链路 A 首批故障场景}
  \label{tab:first-faults}
  \begin{tabularx}{\textwidth}{p{3cm}p{4.5cm}X}
    \toprule
    场景 & DSL & 假设 \\
    \midrule
    geocode-latency & \code{tool.geocoding.latency=4000} & 成功率应基本保持，端到端 p95/p99 上升，Trace 中地理编码 Span 命中延迟 \\
    attraction-exception & \code{rpc.attraction.GetAttractionsNearby.exception=RuntimeError} & HTTP 可能仍成功，但景点候选减少，fallback 与内容质量退化变得可见 \\
    itinerary-unavailable-10p & \code{rpc.itinerary.create.error=UNAVAILABLE,probability=0.1} & 约一部分请求在持久化阶段失败，HTTP 502 上升，成功率下降 \\
    cascade-degrade & \code{tool.geocoding.latency=3000;rpc.attraction.GetAttractionsNearby.exception=RuntimeError} & 延迟与内容质量同时退化，高压下可能放大排队和失败 \\
    \bottomrule
  \end{tabularx}
\end{table}

需要注意，带概率的故障不能只根据 DSL 推断命中次数。论文结果中必须区分“声明概率”和“实际命中”。实际命中应以 Tempo 中 \code{fault.injected=true} 的 Span 数为准。

\section{链路 B 实验设计}

链路 B 的实验重点是多轮状态和控制流鲁棒性。建议设计固定对话脚本，使用户先加载或创建一个 itinerary，再执行新增天数、修改活动、重新生成某天等操作。每轮对话使用相同初始状态和相同用户意图，便于比较 baseline 与 fault 条件。

可选故障包括：

\begin{itemize}
  \item \code{state.itinerary.clear=true}：在状态写回前清空 itinerary，观察 reducer 和前端同步是否阻止空状态覆盖。
  \item \code{state.pending\_day\_plan.clear=true}：清空草案计划，观察用户是否收到可理解提示。
  \item \code{route.should\_continue.force\_route=\_\_end\_\_}：强制提前结束，观察 Agent 是否漏掉应执行工具。
  \item \code{llm.chat\_agent.latency=...} 或 \code{llm.chat\_agent.exception=...}：观察聊天入口是否抛错以及用户侧体验。
\end{itemize}

链路 B 的指标不仅包括 HTTP 成功率和耗时，还应包括 turn 完成率、工具调用次数、状态字段完整性、前端 StateSnapshot 是否更新、用户可见错误类型和 fallback reason。由于当前尚未完成系统化实验，结果表格应保留为待补充。

\section{链路 C 实验设计}

链路 C 的实验重点是跨 Agent 和跨服务可靠性。建议设计订单相关固定任务，例如“查询某 SKU 并创建订单草稿”，并在 baseline 中确认远程 \code{order\_assistant}、商品服务和订单服务均可用。

可选故障包括：

\begin{itemize}
  \item \code{agent.order\_assistant.drop=true}：模拟远程 Agent 不可用，观察主 Agent 是否有明确 fallback。
  \item \code{a2a.trace.drop=true}：模拟 A2A Trace 关联断开，观察 Tempo 中证据链是否不完整。
  \item \code{rpc.product.GetSkuById.error=UNAVAILABLE}：模拟商品服务不可用，观察远端 Agent 的错误处理。
  \item \code{rpc.order.CreateOrder.error=UNAVAILABLE}：模拟订单持久化失败，观察草稿状态和用户提示。
\end{itemize}

链路 C 的关键指标包括远端 Agent 解析成功率、A2A 调用成功率、订单草稿创建率、订单提交成功率、Trace 跨服务完整性和用户可见错误。当前论文不填充结果数值。

\section{指标体系}

\subsection{Locust 指标}

Locust 提供端到端负载指标，适合观察用户视角的退化：

\begin{itemize}
  \item \code{requests/s}：实际吞吐，判断高压组是否达到目标。
  \item \code{failures/s} 和失败率：观察 HTTP 5xx、超时和连接错误。
  \item p50、p95、p99：观察延迟分布，尤其是高延迟故障是否传导到尾延迟。
  \item 平均响应时间：辅助描述整体退化。
  \item 响应体长度：作为内容退化的粗略信号。
\end{itemize}

\subsection{Trace 指标}

Trace 用于确认故障是否实际命中，以及命中后系统走了哪条路径。核心字段包括 \code{experiment.id}、\code{fault.scenario}、\code{fault.injected}、\code{fault.target}、\code{fault.primitive}、\code{fault.outcome}、\code{tool.outcome} 和 \code{tool.fallback\_reason}。

典型 TraceQL 包括：

\begin{verbatim}
{ resource.service.name = "trip-itinerary-planner"
  && .experiment.id = "$experiment_id" }
\end{verbatim}

\begin{verbatim}
{ resource.service.name = "trip-itinerary-planner"
  && .experiment.id = "$experiment_id"
  && .fault.injected = true }
\end{verbatim}

第一条查询用于找到本次实验全部 Trace，第二条用于确认实际注入命中。

\subsection{业务质量指标}

AI 原生应用不能只看 HTTP 2xx。链路 A 至少应记录：

\begin{itemize}
  \item \code{itinerary\_saved\_rate}：HTTP 成功且响应包含持久化 ID 的比例。
  \item \code{day\_plan\_count}：生成行程天数，判断是否退化为空骨架。
  \item \code{activity\_count}：活动数量，判断景点或住宿是否缺失。
  \item \code{markdown\_length}：Markdown 文本长度，粗略衡量生成内容是否异常缩短。
  \item \code{fallback\_rate}：通过 Trace 或响应结构统计 fallback 是否触发。
\end{itemize}

链路 B 和 C 需要扩展相应业务指标，例如 turn 完成率、状态完整性、订单草稿创建率和订单提交率。

\subsection{系统指标}

为了避免把宿主机或容器瓶颈误判为应用鲁棒性问题，Prometheus/Grafana 应补充以下指标：

\begin{itemize}
  \item 宿主机 CPU、内存、磁盘 IO 和网络 IO。
  \item Docker 容器 CPU、内存、网络和重启次数。
  \item Java 服务 JVM、线程池和 gRPC/HTTP 指标。
  \item OTel spanmetrics 转换得到的入口延迟、Span 计数、错误数和故障命中数。
\end{itemize}

这些指标当前仍有一部分尚待实现。论文最终版应在实验环境搭建完成后补充实际看板截图或导出数据。

\section{实验记录模板}

每次实验建议使用 YAML 记录，保证复现条件完整：

\begin{verbatim}
experiment:
  id: "locust-fault-low-geocode-latency-YYYYMMDDHHMM"
  matrix_cell: "fault-low"
  scenario_name: "geocode_latency"
  fault_dsl: "tool.geocoding.latency=4000,probability=1.0"
  target_chain: "A"
  endpoint: "POST /api/v1/itineraries/plannings"
  hypothesis: "低压下高德 4s 延迟会显著提高 p95，但不应明显降低成功率。"
load:
  users: 10
  spawn_rate: 2
  duration: "15m"
windows:
  warmup: "待补充"
  measure: "待补充"
  recovery: "待补充"
locust:
  artifact_dir: "artifacts/locust/..."
  total_requests: "待补充"
  failure_rate: "待补充"
  p95_ms: "待补充"
trace_evidence:
  fault_injected_count: "待补充"
  representative_trace_ids: []
conclusion: "待补充"
\end{verbatim}

该模板的关键是把实验条件和结果分开。配置可以在实验前确定，结果必须在运行后填写。

\section{结果表格占位}

表 \ref{tab:result-placeholder} 是链路 A 四象限实验的结果占位。当前不填数值，等待实际实验后补充。

\begin{table}[htbp]
  \centering
  \caption{链路 A 四象限实验结果占位}
  \label{tab:result-placeholder}
  \begin{tabularx}{\textwidth}{p{2.5cm}XXXX}
    \toprule
    实验 & 请求数 & 失败率 & p95 延迟 & 业务质量结论 \\
    \midrule
    baseline-low & \TODO{待实验} & \TODO{待实验} & \TODO{待实验} & \TODO{待实验} \\
    fault-low & \TODO{待实验} & \TODO{待实验} & \TODO{待实验} & \TODO{待实验} \\
    baseline-high & \TODO{待实验} & \TODO{待实验} & \TODO{待实验} & \TODO{待实验} \\
    fault-high & \TODO{待实验} & \TODO{待实验} & \TODO{待实验} & \TODO{待实验} \\
    \bottomrule
  \end{tabularx}
\end{table}

图 \TODO{待补充 Locust 延迟曲线}、图 \TODO{待补充 Tempo 代表 Trace 截图} 和图 \TODO{待补充 Grafana 资源看板} 应在最终实验后补充。论文最终版需要把每个图与实验结论对应起来，而不是只展示截图。

\section{实验有效性讨论}

为了保证实验结论可信，本文后续实践需要注意以下威胁：

\begin{itemize}
  \item \textbf{LLM 不确定性。} 相同输入可能得到不同输出，业务质量指标应关注结构性变化，并保留请求样本。
  \item \textbf{外部 API 波动。} 高德、模型网关或服务发现本身可能波动，应通过 baseline 和恢复窗口排除背景噪声。
  \item \textbf{机器资源瓶颈。} 本地机器性能不足可能导致高压组失败，必须用系统指标佐证。
  \item \textbf{概率故障统计。} 低概率故障需要足够请求数，否则实际命中数可能不足以支持结论。
  \item \textbf{Trace 采样。} 如果后续启用采样，需要保证故障 Trace 不被丢弃，或采用 tail-based sampling。
\end{itemize}

这些限制不会否定故障注入机制本身，但会影响实验结果的解释。论文最终版应在结果分析中明确说明实验范围和局限。
